\documentclass[12pt]{article}
% Préambule commun à tous les documents extraits de « La Revue CAPES »
\usepackage[T1]{fontenc}
\usepackage[utf8]{inputenc}
\usepackage{lmodern}
\usepackage{amsmath,amssymb,amsthm,amsfonts,mathrsfs}
\usepackage{setspace}
\usepackage{listings}
\usepackage{pgfplots}
\pgfplotsset{compat=1.18}
\usepackage{longtable}
\usepackage{adjustbox}
\usepackage{lettrine}
\usepackage{tkz-tab}
\usepackage{changepage}
\usepackage{pdflscape}
\usepackage{graphicx}
\usepackage{tikz}
\usepackage{xcolor}
\usepackage[a4paper,top=2cm,bottom=2.5cm,left=2.5cm,right=2cm]{geometry}
\usepackage{fancyhdr}
\usepackage{draftwatermark}
\SetWatermarkText{La RevueCapes}
\SetWatermarkLightness{0.9}
\SetWatermarkScale{2}
\usepackage{hyperref}
\hypersetup{colorlinks=true,linkcolor=blue!50!black,urlcolor=blue!50!black}

\definecolor{revue}{RGB}{20,60,120}

\pagestyle{fancy}
\fancyhf{}
\renewcommand{\headrulewidth}{0.4pt}
\setlength{\headheight}{15pt}

% \entete{Matière}{Titre}{Type de document}
\newcommand{\entete}[3]{%
  \fancyhead[L]{\small\textsc{La Revue CAPES} -- #1}%
  \fancyhead[R]{\small #2}%
  \fancyfoot[C]{\thepage}%
  \thispagestyle{plain}%
  \begin{center}
    {\color{revue}\rule{\textwidth}{1.2pt}}\\[4pt]
    {\small\textsc{Concours directs CAP-CEG / CAPES -- Burkina Faso}}\\[6pt]
    {\LARGE\bfseries\color{revue} #2}\\[6pt]
    {\large\itshape #1 -- #3}\\[2pt]
    {\color{revue}\rule{\textwidth}{1.2pt}}
  \end{center}
  \vspace{0.5cm}%
}

% Encadré pour les remarques ajoutées dans les corrigés
\newcommand{\remarque}[1]{\par\medskip\noindent\fcolorbox{revue}{revue!6}{\parbox{\dimexpr\linewidth-2\fboxsep-2\fboxrule}{\textbf{\color{revue}Remarque.} #1}}\par\medskip}


\begin{document}
\entete{Culture générale}{Corrigé du sujet type de dissertation n°10}{Correction}
\begin{spacing}{1.5}

\noindent\textbf{Sujet.} L'accès à l'éducation pour tous au Burkina Faso : enjeux et défis.

\rubrique{1. Analyse du sujet}

\begin{itemize}
\item \textbf{Type de sujet :} sujet-thème, sans question explicite. La formule « enjeux et défis » suggère un \textbf{plan analytique} : I. les enjeux (pourquoi l'éducation pour tous est importante), II. les défis (les obstacles), III. les solutions (comment relever ces défis).
\item \textbf{Mots-clés :}
  \begin{itemize}
  \item « accès à l'éducation pour tous » : possibilité pour chaque enfant, chaque jeune et chaque adulte, sans distinction de sexe, de région, de revenu ou de handicap, de bénéficier d'une éducation de qualité ; c'est un droit reconnu par la Constitution et par les conventions internationales, et l'un des Objectifs de développement durable (ODD 4) ;
  \item « enjeux » : ce que l'on peut gagner ou perdre ;
  \item « défis » : les difficultés à surmonter.
  \end{itemize}
\item \textbf{Problématique :} pourquoi l'éducation pour tous est-elle un objectif essentiel pour le Burkina Faso, et quels obstacles faut-il surmonter pour l'atteindre ?
\end{itemize}

\rubrique{2. Plan détaillé (plan analytique)}

\begin{enumerate}
\item[I.] \textbf{Les enjeux.} 1. Un enjeu de droit et d'égalité. 2. Un enjeu économique : capital humain et développement. 3. Un enjeu de paix, de citoyenneté et de cohésion sociale.
\item[II.] \textbf{Les défis.} 1. L'insuffisance des infrastructures, des enseignants et des moyens. 2. Les inégalités (filles, zones rurales, pauvres, handicap). 3. L'insécurité et les déplacements de population.
\item[III.] \textbf{Les pistes pour relever ces défis.} 1. Investir davantage et mieux. 2. Lever les obstacles sociaux et économiques. 3. Innover : éducation en situation d'urgence, éducation non formelle, numérique.
\end{enumerate}

\rubrique{3. Proposition de rédaction}

\partie{Introduction}

« L'éducation est l'arme la plus puissante qu'on puisse utiliser pour changer le monde », affirmait Nelson Mandela. Consciente de cette vérité, la communauté internationale a fait de l'éducation pour tous l'un de ses objectifs prioritaires, repris dans l'Objectif de développement durable n°4 : « assurer l'accès de tous à une éducation de qualité ». Au Burkina Faso, la loi d'orientation de l'éducation de 2007 a rendu l'enseignement de base obligatoire de six à seize ans. D'importants progrès ont été réalisés depuis les années 2000, mais de nombreux enfants restent exclus de l'école ou la quittent prématurément. Quels sont alors les enjeux de l'accès à l'éducation pour tous au Burkina Faso, et quels défis faut-il relever pour l'atteindre ? Nous examinerons d'abord les enjeux de l'éducation pour tous, puis les défis qui en freinent la réalisation, avant de proposer des pistes de solution.

\partie{I. Les enjeux de l'accès à l'éducation pour tous}

\souspartie{1. Un enjeu de droit et d'égalité}
L'éducation est un droit fondamental reconnu par la Déclaration universelle des droits de l'homme (article 26), la Convention relative aux droits de l'enfant et la Constitution burkinabè. Garantir l'accès de tous à l'école, c'est donner à chaque enfant, qu'il soit fille ou garçon, riche ou pauvre, citadin ou villageois, les mêmes chances de réussir dans la vie.

\souspartie{2. Un enjeu économique}
Aucun pays ne s'est développé sans une population instruite. L'éducation forme les agriculteurs, techniciens, enseignants, médecins et entrepreneurs dont le pays a besoin ; elle augmente la productivité et les revenus, et contribue à réduire la pauvreté. Une population alphabétisée adopte plus facilement les innovations agricoles et sanitaires.

\souspartie{3. Un enjeu de paix et de citoyenneté}
L'école enseigne les valeurs de tolérance, de respect et de vivre-ensemble ; elle développe l'esprit critique. Une jeunesse instruite est moins vulnérable aux discours extrémistes et à l'enrôlement par les groupes armés. L'éducation est donc aussi un rempart pour la paix et la cohésion nationale.

\transition{Les enjeux sont donc considérables. Pourtant, l'éducation pour tous reste un objectif difficile à atteindre au Burkina Faso.}

\partie{II. Les défis de l'éducation pour tous}

\souspartie{1. L'insuffisance des moyens}
Le nombre d'écoles, de salles de classe, d'enseignants qualifiés et de manuels reste insuffisant face à une population scolaire en forte croissance. Classes surchargées, écoles sous paillote, éloignement des établissements, manque de cantines et de points d'eau : ces conditions découragent les familles et nuisent à la qualité de l'enseignement.

\souspartie{2. Les inégalités d'accès}
Les enfants des zones rurales, des familles pauvres, les filles et les enfants en situation de handicap sont les plus touchés par la non-scolarisation et l'abandon. La pauvreté pousse certains enfants vers le travail (agriculture, élevage, commerce, orpaillage). Les filles sont freinées par les travaux domestiques, les mariages et grossesses précoces et certains préjugés.

\souspartie{3. L'insécurité}
Depuis 2015, les attaques des groupes armés ont entraîné la fermeture de milliers d'écoles et privé plusieurs centaines de milliers d'élèves de scolarité. Les enseignants sont menacés, les familles déplacées, les écoles des zones d'accueil surchargées. L'insécurité est aujourd'hui le défi le plus grave pour l'éducation au Burkina Faso.

\transition{Face à ces défis, plusieurs pistes peuvent être envisagées pour rapprocher le pays de l'objectif de l'éducation pour tous.}

\partie{III. Des pistes pour relever ces défis}

\souspartie{1. Investir davantage et mieux}
Il s'agit de construire et d'équiper des écoles, de recruter et de former des enseignants, de généraliser les cantines scolaires et les kits scolaires, et de lutter contre la mauvaise gestion des ressources. La décentralisation de la gestion de l'éducation de base aux communes doit être accompagnée de moyens suffisants.

\souspartie{2. Lever les obstacles sociaux et économiques}
Sensibiliser les parents à l'importance de l'école, en particulier pour les filles ; lutter contre les mariages précoces ; accorder des bourses et des aides aux familles vulnérables ; rendre les écoles accessibles aux enfants handicapés. L'implication des associations de parents d'élèves, des mères éducatrices et des leaders communautaires est essentielle.

\souspartie{3. Innover face aux crises}
L'éducation en situation d'urgence permet de maintenir l'école dans les zones touchées par l'insécurité : écoles temporaires, enseignement à distance par la radio ou le numérique, accueil des élèves déplacés. L'éducation non formelle (centres d'alphabétisation, écoles bilingues, formation professionnelle) offre une seconde chance aux enfants et aux adultes qui n'ont pas pu aller à l'école.

\partie{Conclusion}

L'accès à l'éducation pour tous est un enjeu majeur pour le Burkina Faso : il conditionne l'égalité des chances, le développement économique et la paix. Mais cet objectif se heurte à l'insuffisance des moyens, aux inégalités d'accès et, surtout, à l'insécurité qui a fermé de nombreuses écoles. Pour le relever, il faut investir davantage et mieux, lever les obstacles sociaux et innover pour maintenir l'école en temps de crise. L'éducation pour tous n'est pas seulement l'affaire de l'État : c'est une responsabilité partagée par les familles, les communautés et l'ensemble de la nation.

\end{spacing}
\end{document}
